% 02-introduction.tex
\chapter{Introduction}
\label{ch:introduction}

\section{Background}
\label{sec:background}

Container technology has transformed how software is developed,
deployed, and operated. Since the introduction of Docker in 2013,
containers have become the standard unit of packaging applications,
providing isolation, reproducibility, and portability across computing
environments. The container ecosystem has expanded to include
orchestration platforms~\cite{burns2016borg}, image registries,
runtime specifications~\cite{open2021runtime}, and networking
primitives that collectively enable modern cloud-native architectures.

However, this ecosystem was designed for server-class hardware running
Linux distributions with full kernel access. Mobile devices, which
represent the majority of computing hardware deployed worldwide, have
been largely excluded from the container revolution. Android, which
powers over 3 billion active devices~\cite{statcounter2025}, presents
fundamental barriers to containerization that existing tools cannot
address.

\section{The Mobile Containerization Problem}
\label{sec:mobile-problem}

Android devices face several constraints that prevent traditional
container runtimes from functioning:

\begin{enumerate}[leftmargin=*, itemsep=4pt]
  \item \textbf{Restricted syscalls.} The Android kernel disables or
        restricts numerous system calls required by container runtimes,
        including \texttt{clone()} with certain namespace flags,
        \texttt{pivot\_root()}, and \texttt{mount()} for bind mounts.
        SELinux Mandatory Access Control (MAC) policies further limit
        what operations are permitted.

  \item \textbf{No root access.} Modern Android devices use verified
        boot chains and dm-verity to prevent unauthorized modifications.
        Gaining root access typically requires unlocking the bootloader,
        which voids warranties and may compromise device security.

  \item \textbf{Networking restrictions.} The \texttt{iptables}
        subsystem is unavailable without root on most Android devices.
        Port 53 (DNS) is blocked by SELinux policies. The network
        stack cannot be manipulated through standard Linux networking
        tools.

  \item \textbf{Filesystem limitations.} Android uses ext4 or F2FS
        with encrypted user partitions. Standard overlay filesystem
        operations require kernel support that may not be available,
        and the storage layout differs significantly from desktop
        Linux distributions.

  \item \textbf{Process management.} Android's Zygote process model
        and application lifecycle management interfere with long-running
        container processes. Background process restrictions can terminate
        container runtimes unexpectedly.
\end{enumerate}

These constraints mean that Docker, containerd, and other
server-oriented runtimes cannot operate on Android without significant
modifications. The result is a gap in the tooling ecosystem that
limits what developers can do on mobile devices.

\section{The Doki Approach}
\label{sec:approach}

Doki addresses these challenges through a combination of techniques
that work within Android's constraints rather than against them:

\begin{itemize}[leftmargin=*, itemsep=4pt]
  \item \textbf{Proot-based syscall interception.} Instead of relying
        on kernel-level namespaces, Doki uses
        \texttt{proot}~\cite{proot2024} to intercept system calls via
        \texttt{ptrace()}. This provides a virtualized filesystem view
        and process namespace without requiring root access or kernel
        modifications.

  \item \textbf{Filesystem detection probes.} Doki detects its
        execution environment by probing the filesystem (checking for
        \texttt{/data/data/com.termux}) rather than relying on
        \texttt{runtime.GOOS}, which returns \texttt{linux} on both
        Android and standard Linux. This enables environment-specific
        behavior without compile-time restrictions.

  \item \textbf{Adaptive networking.} When \texttt{iptables} is
        unavailable, Doki falls back to alternative port forwarding
        mechanisms. DNS operates on a non-standard port (8053) to
        avoid SELinux restrictions. The networking stack auto-detects
        available capabilities and adjusts accordingly.

  \item \textbf{Static compilation.} All Doki binaries are compiled
        as statically linked executables with \texttt{CGO\_ENABLED=0}.
        This eliminates runtime dependencies on shared libraries that
        may not be present on Android devices.

  \item \textbf{Multi-runtime architecture.} Doki supports 12
        isolation levels, from WASM sandboxes to hardware-level
        microVMs. The runtime registry probes the host at startup and
        automatically selects the strongest isolation mode available,
        providing graceful degradation across different device
        capabilities.
\end{itemize}

\section{Design Principles}
\label{sec:principles}

Doki's design is guided by the following principles:

\begin{description}[leftmargin=*, style=nextline, font=\sffamily\bfseries]
  \item[Rootless by default.]
    Every feature must work without root access. Root is never assumed,
    requested, or required for core functionality.

  \item[Progressive enhancement.]
    Doki detects available capabilities at startup and enables features
    accordingly. A device with limited kernel support still gets a
    functional container runtime, just with weaker isolation.

  \item[Docker compatibility.]
    Where possible, Doki accepts Dockerfile syntax, Docker Compose
    files, and OCI image formats. This lowers the barrier to adoption
    for developers already familiar with the Docker ecosystem.

  \item[Single-binary distribution.]
    The entire system is distributed as four statically linked
    binaries totaling approximately 28.3\,MB. No package manager,
    runtime, or system-level installation is required beyond
    extracting the archive.

  \item[Security in depth.]
    Isolation is achieved through multiple layered mechanisms:
        syscall filtering (seccomp), mandatory access control
        (AppArmor), Linux capabilities, and user namespaces.
        Each layer provides independent protection even if another
        layer is bypassed.
\end{description}

\section{Contributions}
\label{sec:contributions}

This report makes the following contributions:

\begin{enumerate}[leftmargin=*, itemsep=4pt]
  \item A comprehensive architecture for rootless containerization
        on Android, demonstrating that Docker-like functionality is
        achievable without root access or kernel modifications.

  \item A 12-level isolation framework that automatically selects the
        strongest available runtime based on host capabilities,
        providing both security and compatibility.

  \item DokiLink-Lite, a peer-to-peer mesh networking protocol with
        TLS~1.3 encryption and optional NaCl secretbox payload
        encryption for secure inter-device communication.

  \item An empirical evaluation of container startup times, memory
        usage, and network throughput across multiple isolation levels
        and device configurations.
\end{enumerate}

\section{Document Organization}
\label{sec:organization}

The remainder of this report is organized as follows.
\cref{ch:architecture} presents the system architecture, including
the daemon design, CLI structure, and communication protocols.
\cref{ch:components} describes the individual subsystems in detail.
\cref{ch:isolation} covers the 12 isolation levels and their
implementations. \cref{ch:networking} explains the networking stack,
including bridge networking, internal DNS, and DokiLink-Lite mesh.
\cref{ch:storage} describes image storage, layer management, and
the overlay filesystem. \cref{ch:build} covers the build system,
Dockerfile support, and Compose orchestration. \cref{ch:metrics}
presents performance benchmarks and comparisons. \cref{ch:roadmap}
outlines future development plans.
